% Responsable: QUINTERO RICO LUIS DAVID.
\section{Plan del proyecto}
\label{sec:plan-proyecto}

\subsection{Roles y responsabilidades}
El equipo tiene tres integrantes y cuatro roles; la Líder de Desarrollo asume también la responsabilidad de calidad. La Tabla~\ref{tab:roles} resume las responsabilidades de cada uno.

\begin{table}[htbp]
\centering\small
\caption{Roles y responsabilidades del equipo}
\label{tab:roles}
\begin{tabularx}{\textwidth}{@{}>{\raggedright\arraybackslash}p{3.6cm}>{\raggedright\arraybackslash}p{3.3cm}>{\raggedright\arraybackslash}X@{}}
\toprule
\textbf{Integrante} & \textbf{Rol} & \textbf{Responsabilidades}\\
\midrule
SÁNCHEZ VERA DARWIN JOEL & Arquitecto & Estructura del repositorio, trabajos relacionados, ADR-001 y ADR-002, diagramas C4, modelo de comunicación, CAP/PACELC, contratos de I-05 e I-06, servicio B de postulaciones y riesgos derivados de FORM-AC-01.\\
MUÑOZ VERA MELANIE MICHELLE & Líder de Desarrollo y Responsable de Calidad & Reglas de trabajo del repositorio, servicio A de elegibilidad y SGA simulado, prueba de comunicación entre servicios, RNF, análisis de riesgos, matriz de trazabilidad, esqueleto funcional y capturas de ejecución.\\
QUINTERO RICO LUIS DAVID & Responsable de Documentación & Base LaTeX del documento, guion y acta de la entrevista, contexto, objetivos, requisitos funcionales, referencias de contexto, plan del proyecto, tabla de aportes, README y bitácora.\\
\bottomrule
\end{tabularx}
\end{table}

\subsection{Tablero de seguimiento}
El trabajo se gestiona en el tablero de GitHub Projects «SGAC -- Plan del proyecto (E1 a E4)», disponible en \url{https://github.com/users/DarwinDevCode/projects/14}, vinculado al repositorio \url{https://github.com/DarwinDevCode/AD_TA-PFC-E1}. El tablero presenta los issues en una vista de tabla con los siguientes campos, que permiten agrupar y filtrar el trabajo:
\begin{itemize}
  \item \textbf{Estado:} Por hacer, En curso, En revisión y Hecho.
  \item \textbf{Entrega:} E1, E2, E3 y E4.
  \item \textbf{Componente:} documento, arquitectura, servicio B, servicio A, SGA simulado, contratos de API, mensajería, servicios de evaluación, seguimiento y notificaciones, contenedores, evidencias y gestión.
  \item \textbf{Rol:} Arquitecto, Desarrollo y calidad, y Documentación.
\end{itemize}
Cada entrega tiene un issue principal con su alcance y fechas (\#9 para E1, \#8 para E2, \#12 para E3 y \#13 para E4), y los cambios se integran mediante pull requests.

\subsection{Issues de E1 y E2 por integrante}
La Tabla~\ref{tab:issues} agrupa los issues registrados en el repositorio según su responsable. El issue \#9 de E1 está asignado a los tres integrantes.

\begingroup\small
\begin{longtable}{@{}>{\raggedright\arraybackslash}p{.24\textwidth}>{\raggedright\arraybackslash}p{.08\textwidth}>{\raggedright\arraybackslash}p{.48\textwidth}>{\raggedright\arraybackslash}p{.11\textwidth}@{}}
\caption{Issues de E1 y E2 por integrante}\label{tab:issues}\\
\toprule
\textbf{Integrante} & \textbf{Issue} & \textbf{Título} & \textbf{Entrega}\\
\midrule\endfirsthead
\toprule\textbf{Integrante} & \textbf{Issue} & \textbf{Título} & \textbf{Entrega}\\\midrule\endhead
\bottomrule\endfoot
SÁNCHEZ VERA DARWIN JOEL & \#14 & Preparar la estructura de carpetas del repositorio & E1\\
 & \#15 & Redactar los trabajos relacionados y las referencias de arquitectura & E1\\
 & \#16 & Elaborar el diagrama C4 de nivel 1 & E1\\
 & \#17 & Registrar ADR-001 sobre el estilo de arquitectura & E1\\
 & \#19 & Registrar ADR-002 sobre el mecanismo de comunicación & E1\\
 & \#20 & Elaborar el diagrama C4 de nivel 2 & E1\\
 & \#21 & Redactar las secciones 07, 09, 10 y 11 & E1\\
 & \#22 & Definir los contratos de API de I-05 e I-06 & E1\\
 & \#23 & Implementar el servicio B de postulaciones y sus pruebas & E1\\
 & \#29 & Probar el servicio B con el servicio A real & E1\\
 & \#30 & Aportar los riesgos del incidente de FORM-AC-01 & E1\\
 & \#31 & Registrar las fotografías de las sesiones de diseño y el primer bosquejo & E1\\
 & \#7 & Alinear C4, ADR, comunicación y CAP con los requisitos & E1\\
 & \#8 & Microservicios, API y contenedores & E2\\
\addlinespace
MUÑOZ VERA MELANIE MICHELLE & \#1 & Preparar el repositorio y la documentación de ADR & E1\\
 & \#4 & Preparar los borradores de requisitos no funcionales, riesgos y referencias & E1\\
 & \#24 & Implementar el servicio A de elegibilidad y el SGA simulado & E1\\
 & \#25 & Automatizar la prueba de comunicación entre servicios & E1\\
 & \#26 & Redactar la sección del esqueleto funcional y las instrucciones de ejecución & E1\\
 & \#27 & Revisar la coherencia con la matriz de trazabilidad & E1\\
 & \#28 & Registrar las capturas de ejecución del esqueleto & E1\\
\addlinespace
QUINTERO RICO LUIS DAVID & \#6 & Completar la trazabilidad de los RF y la bibliografía de contexto & E1\\
 & \#32 & Crear la base del documento LaTeX & E1\\
\end{longtable}
\endgroup

\subsection{Cronograma de entregas}
La Tabla~\ref{tab:cronograma} presenta las cuatro entregas del Proyecto Fin de Curso. Las semanas se cuentan desde el inicio del periodo académico, el 26 de octubre de 2026.

\begin{table}[htbp]
\centering\small
\caption{Cronograma de las entregas del Proyecto Fin de Curso}
\label{tab:cronograma}
\begin{tabularx}{\textwidth}{@{}p{1.7cm}p{1.7cm}>{\raggedright\arraybackslash}p{3.0cm}>{\raggedright\arraybackslash}X@{}}
\toprule
\textbf{Entrega} & \textbf{Semana} & \textbf{Fechas} & \textbf{Contenido}\\
\midrule
E1 & 5 & 23 al 29-11-2026; entrega el 29-11-2026, 23:55 & Problema, objetivos, RF y RNF, ADR-001 y ADR-002, C4 de niveles 1 y 2, CAP/PACELC, riesgos y esqueleto con los servicios de postulaciones y elegibilidad.\\
\addlinespace
E2 & 9 & 21 al 27-12-2026 & Servicios de ADR-001 con sus API documentadas, contenedores por servicio e intermediario de mensajes para los avisos.\\
\addlinespace
E3 & 13 & 18 al 24-01-2027 & Base de datos propia por servicio, replicación y procesamiento distribuido de los datos del concurso.\\
\addlinespace
E4 & 17 & 15 al 21-02-2027 & Sistema final con aplicación web, pruebas de los RNF, evaluación de calidad y defensa.\\
\bottomrule
\end{tabularx}
\end{table}

\subsection{Actividades de la semana de E1}
\begin{itemize}
  \item 23 y 24 de noviembre: integrar las secciones del documento y alinear los RF con el acta de la entrevista.
  \item 25 y 26 de noviembre: revisar RNF, riesgos y trazabilidad; compilar el informe.
  \item 27 y 28 de noviembre: ejecutar el esqueleto desde un clon limpio, registrar las capturas y grabar el video.
  \item 29 de noviembre: revisar el PDF final y subir la entrega a la plataforma de la asignatura.
\end{itemize}
